\section{Introduction}

Evidence-based fact-checking is a pipeline: a system retrieves source material, decides which items are relevant, and predicts a verdict that should be attributable to those items \cite{zeng2021automated,zhou2020survey}. Multimodal claims make this harder because decisive support may be distributed across text and images. A correct verdict is therefore not sufficient. The evidence set must also be identifiable, reasonably complete, and free of irrelevant material.

Most retrieval pipelines expose a fixed top-$k$ set. This is reproducible, but it assumes the same evidence budget for every claim. Adaptive selection appears to address that mismatch by estimating item relevance and choosing a claim-specific set size. Its evaluation is easily confounded, however. If a released judgment packet already contains most gold evidence, a selector is tested mainly on set decoding. If candidates are retrieved label-blind from a full corpus, missing gold items remain false negatives and high-scoring distractors change both calibration and the precision--recall trade-off. A selector conclusion may therefore reverse without any change to the decoder.

An earlier development study in this project demonstrated exactly this reversal. Calibrated Expected-F1 Evidence Selection (CEFES) beat fixed selection inside a restricted judgment packet but lost under full-corpus retrieval. The immediate question is not simply whether a more complex selector can recover a positive number. It is when an adaptive conclusion is scientifically valid, how retrieval and scoring affect it, and whether any gain survives strong fixed and threshold controls.

We address this with a staged retrieval-conditioned framework. Retrieval is evaluated before a new scorer, scoring before a new decoder, and evidence selection before a shared-verifier safeguard. The proposed \method{} system adds: (i) BGE-M3/BM25 reciprocal-rank fusion and cross-encoder text reranking; (ii) SigLIP2/filename-text image fusion; (iii) modality-specific or cross-encoder-augmented logistic membership scoring; (iv) training-only temperature calibration; and (v) exact expected-F1 decoding with set-size and same-source redundancy penalties. All selectors receive identical packets, features, supervision, and tuning roles.

The central research questions are:
\begin{enumerate}
    \item Does stronger label-blind retrieval change the measured adaptive-versus-fixed conclusion?
    \item Do modality-aware scoring and precision penalties make expected-F1 decoding competitive with fixed, threshold, and count controls?
    \item Under what controlled retrieval, duplication, modality, and calibration changes does the conclusion move?
    \item Does selected evidence preserve verdict macro-F1, reproduce cleanly, and admit an unchanged external replication?
\end{enumerate}

The contributions are:
\begin{itemize}
    \item a retrieval-conditioned formulation and sequential gate that separates missing evidence, within-packet ranking, probability calibration, set decoding, and verifier behavior;
    \item \method{}, a matched multimodal scorer and precision-aware expected-F1 decoder that reduces CEFES over-selection without changing corpus access;
    \item a five-seed, complete-gold MOCHEG study with fixed, threshold, count, original CEFES, size-penalized, and redundancy-penalized controls, paired uncertainty, Holm correction, and a shared-verifier non-inferiority safeguard;
    \item controlled perturbations showing that an apparent adaptive advantage emerges under deliberate top-rank distractor poisoning, while ordinary retrieval and calibration conditions retain a fixed-selector advantage; and
    \item an auditable negative decision: six independent audits pass, 27 core artifacts reproduce exactly, the formal test remains sealed, and an external audit documents why current candidates do not support unchanged confirmation.
\end{itemize}

The paper consequently does not claim that \method{} is superior. It establishes that improved image retrieval and scoring are insufficient when text coverage, ranking, and downstream safeguards do not jointly pass.
